\chapter{Union--Find, Spanning Trees and Topological Order}\label{ch:dsu}

This chapter is about three ideas that organise graphs globally rather than by searching them. The disjoint set union
structure maintains a partition of the vertices under merging, in almost constant time per operation; it answers ``are these
two connected?'' as edges arrive, and with small extensions it maintains equations and parities along the way. Minimum
spanning trees connect everything as cheaply as possible, and their theory (the cut and cycle properties) explains a whole
family of greedy algorithms. Prüfer codes turn labelled trees into sequences and back, which makes counting trees a matter of
counting sequences. And topological order arranges a directed acyclic graph so that every edge points forward, the
precondition for every dynamic programme on a DAG.

\begin{chapterintro}{\gold\ \plat\ \ioi}{Chapter~\ref{ch:graphs} (graph basics) and Section~\ref{sec:greedy} (exchange
arguments).}
  \item implement a DSU with union by size and path compression, and extend it with sizes, offsets and parities;
  \item prove the cut and cycle properties, and use them to justify Kruskal and Prim;
  \item answer offline ``connected using edges of weight at most $x$'' queries, and build Kruskal reconstruction trees;
  \item encode and decode Prüfer sequences, and count labelled trees with given degrees;
  \item compute topological orders, detect cycles with them, and run DP along them.
\end{chapterintro}

% ======================================================================================================================
\section[Disjoint set union]{Disjoint set union \hfill \gold}\label{sec:dsu}

A \defn{disjoint set union} (DSU, also \emph{union--find}) maintains a partition of $\{0, \ldots, n - 1\}$ into disjoint sets
under two operations: \emph{unite} the sets containing $a$ and $b$, and \emph{find} a representative of the set containing
$x$, so that ``are $a$ and $b$ in the same set?'' becomes ``do they have the same representative?''. Each set is stored as a
rooted tree of parent pointers whose root is the representative; uniting two sets links one root under the other
(Figure~\ref{fig:dsuunion}).

\begin{figure}[htbp]
\centering
\begin{tikzpicture}
  \begin{scope}
    \node[vtxs, hot] (r1) at (0, 0) {0};
    \node[vtxs] (a1) at (-0.6, -1) {1}; \node[vtxs] (a2) at (0.6, -1) {2}; \node[vtxs] (a3) at (0.6, -2) {3};
    \draw[ded] (a1) -- (r1); \draw[ded] (a2) -- (r1); \draw[ded] (a3) -- (a2);
    \node[vtxs, hot] (r2) at (2.0, 0) {4}; \node[vtxs] (b1) at (2.0, -1) {5};
    \draw[ded] (b1) -- (r2);
    \node[lbl] at (1.0, -2.75) {two sets, sizes 4 and 2};
  \end{scope}
  \draw[->, line width=1.2pt, draw=fgblue] (3.0, -1) -- node[above, font=\sffamily\scriptsize, text=fgblue]
    {unite(3, 5)} (4.6, -1);
  \begin{scope}[xshift=6.4cm]
    \node[vtxs, hot] (r) at (0, 0) {0};
    \node[vtxs] (c1) at (-1.0, -1) {1}; \node[vtxs] (c2) at (0, -1) {2}; \node[vtxs] (c3) at (0, -2) {3};
    \node[vtxs, cool] (c4) at (1.0, -1) {4}; \node[vtxs] (c5) at (1.0, -2) {5};
    \draw[ded] (c1) -- (r); \draw[ded] (c2) -- (r); \draw[ded] (c3) -- (c2); \draw[bded] (c4) -- (r);
    \draw[ded] (c5) -- (c4);
    \node[lbl, align=center] at (0, -2.85) {the smaller tree's root (4) goes\\under the larger tree's root (0)};
  \end{scope}
\end{tikzpicture}
\caption{A DSU stores each set as a tree of parent pointers (arrows point to the parent); roots (red) are the representatives.
Union by size links the root of the smaller tree under the root of the larger one.}
\label{fig:dsuunion}
\end{figure}

Two optimisations make the operations fast.

\begin{itemize}
  \item \textbf{Union by size}: when uniting, link the root of the smaller tree under the root of the larger one. A vertex
        then gets deeper only when its tree at least doubles in size, which can happen at most $\log_2 n$ times, so every
        tree has depth $\Oh{\log n}$.
  \item \textbf{Path compression}: during \emph{find}, make the vertices on the path point directly to the root, or at least
        closer to it (Figure~\ref{fig:pathcomp}). Later finds on the same path become $\Oh{1}$.
\end{itemize}

\noindent With both, a sequence of operations costs $\Oh{\alpha(n)}$ amortised per operation, where $\alpha$ is the inverse
Ackermann function, which is at most 4 for any $n$ that fits in the universe. For contest purposes, DSU operations are
constant time. The template uses \emph{path halving} (\cc{x = par[x] = par[par[x]]}), which gets the same bound with a loop
instead of recursion.

\begin{figure}[htbp]
\centering
\begin{tikzpicture}
  \begin{scope}
    \foreach \vv/\yy in {0/0, 1/-0.9, 2/-1.8, 3/-2.7, 4/-3.6} \node[vtxs] (p\vv) at (0, \yy) {\vv};
    \node[vtxs, hot] at (p0) {0};
    \node[vtxs, cool] at (p4) {4};
    \foreach \ea/\eb in {1/0, 2/1, 3/2, 4/3} \draw[ded] (p\ea) -- (p\eb);
    \node[lbl] at (0, -4.35) {before find(4)};
  \end{scope}
  \draw[->, line width=1.2pt, draw=fgblue] (1.1, -1.8) -- (2.4, -1.8);
  \begin{scope}[xshift=4.3cm]
    \node[vtxs, hot] (q0) at (0, 0) {0};
    \foreach \vv/\xx in {1/-1.8, 2/-0.6, 3/0.6, 4/1.8} \node[vtxs] (q\vv) at (\xx, -1.3) {\vv};
    \node[vtxs, cool] at (q4) {4};
    \foreach \a in {1, 2, 3, 4} \draw[ded] (q\a) -- (q0);
    \node[lbl, align=center] at (0, -2.2) {after full compression:\\every vertex on the path points to the root};
  \end{scope}
\end{tikzpicture}
\caption{Path compression. A \emph{find} on the deepest vertex walks the whole path once, and afterwards every vertex on it is
one step from the root.}
\label{fig:pathcomp}
\end{figure}

\usetemplate{DSU (union-find)}{Ctrl+Alt+T, U}{cpdsu}

The complete use of a DSU for incremental connectivity, printing the number of components and the largest component after
every new road, is the worked problem Road Construction (Chapter~\ref{ch:worked}). When the edges are \emph{removed} instead,
reverse time so that they are added (Section~\ref{sec:reverse}).

\begin{recognize}
Merging groups and asking whether two things are in the same group; connectivity while edges are only added; Kruskal;
``the first moment two things become connected'' (add edges by time or weight); equivalence classes (``$a$ and $b$ have the
same colour''); offline queries sorted by a threshold.
\end{recognize}

\subsection{Beyond connectivity: sizes, offsets and parities}

The root of a set is a natural place to store information about the whole set: its size, its sum, its minimum, a list of
members. Uniting merges the information of the two roots. More interesting is information stored on the \emph{edges} of the
DSU trees: each vertex stores a value relative to its parent, and a find accumulates these values along the path to the root.

\begin{problem}{Consistent differences}
There are $n$ unknowns. Process $q$ statements ``$x_b - x_a = c$'', reporting whether each one is consistent with the accepted
statements so far (and accepting it if so), and queries ``what is $x_b - x_a$?'', answering \texttt{?} if it is not determined.
\solution A \defn[weighted DSU]{weighted DSU}. Every vertex $v$ stores $\mathit{off}(v) = x_v - x_{\mathit{par}(v)}$; a find
with path compression sums the offsets on the way, so afterwards $\mathit{off}(v) = x_v - x_{\mathit{root}}$
(Figure~\ref{fig:weighteddsu}). Two vertices in the same set have a determined difference $\mathit{off}(b) - \mathit{off}(a)$.
Uniting two sets for a new statement sets the offset of one root relative to the other so that the statement holds.
\end{problem}

\begin{figure}[htbp]
\centering
\begin{tikzpicture}
  \node[vtx, hot] (r) at (0, 0) {1};
  \node[vtx] (v2) at (-1.5, -1.4) {2};
  \node[vtx] (v3) at (1.5, -1.4) {3};
  \node[vtx] (v4) at (1.5, -2.9) {4};
  \draw[ded] (v2) -- node[wl] {$+3$} (r);
  \draw[ded] (v3) -- node[wl] {$-1$} (r);
  \draw[ded] (v4) -- node[wl] {$+2$} (v3);
  \node[notel, anchor=west, text width=6.4cm] at (3.0, -1.2) {Edge labels are $\mathit{off}(v) = x_v - x_{\mathit{par}(v)}$.
    Relative to the root: $x_2 - x_1 = 3$, $x_3 - x_1 = -1$, $x_4 - x_1 = -1 + 2 = 1$. So the query $x_4 - x_2$ answers
    $1 - 3 = -2$, and the statement $x_4 - x_2 = 5$ is a conflict.};
\end{tikzpicture}
\caption{A weighted DSU: offsets along the path to the root add up to the difference between a vertex and its root.}
\label{fig:weighteddsu}
\end{figure}

\begin{code}
void solve() {
    int n, q;
    cin >> n >> q;
    // par[v] and off[v] = x_v - x_par[v]; the root r has off[r] = 0
    vi par(n + 1), off(n + 1, 0), siz(n + 1, 1);
    iota(all(par), 0);
    // find also makes off[v] = x_v - x_root (path compression)
    auto find = [&](auto&& self, int v) -> int {
        if (par[v] == v) return v;
        int r = self(self, par[v]);
        off[v] += off[par[v]];  // par[v] now points at r directly
        par[v] = r;
        return r;
    };
    while (q--) {
        int type, a, b;
        cin >> type >> a >> b;
        int ra = find(find, a), rb = find(find, b);
        if (type == 1) {  // the fact x_b - x_a = c
            ll c;
            cin >> c;
            if (ra == rb) {
                print(off[b] - off[a] == c ? "OK" : "CONFLICT");
                continue;
            }
            // link the roots so that x_b - x_a = c holds
            ll d = off[a] + c - off[b];  // the required x_rb - x_ra
            if (siz[ra] < siz[rb]) swap(ra, rb), d = -d;
            par[rb] = ra, off[rb] = d, siz[ra] += siz[rb];
            print("OK");
        } else {  // query x_b - x_a
            if (ra != rb) print("?");
            else print(off[b] - off[a]);
        }
    }
}
\end{code}

The same structure with offsets modulo 2 is the \emph{parity DSU}: it maintains constraints ``$a$ and $b$ are in different
groups'' and ``in the same group'', detects the first contradiction (an odd cycle), and so checks bipartiteness online as edges
arrive. With offsets in any group (addition modulo $k$, XOR, permutations composed along the path) the idea is the same.

\begin{trick}[What a DSU cannot do, and the workarounds]
A DSU cannot split a set. If edges are only deleted, reverse time so that they are inserted. If insertions and deletions mix,
answer the queries offline with a \emph{DSU with rollback} (union by size without path compression, keeping a stack of changes
to undo) over a segment tree on the time axis, in $\Oh{m \log m \log n}$ (Chapter~\ref{ch:connectivity}). And when sets carry
lists that must really be merged, merge the smaller list into the larger one: every element moves $\Oh{\log n}$ times
(small to large, Chapter~\ref{ch:trees}).
\end{trick}

\usetemplate{DSU with rollback}{Ctrl+Alt+Y, R}{cpdsurb}

% ======================================================================================================================
\section[Minimum spanning trees]{Minimum spanning trees \hfill \gold}\label{sec:mst}

A \defn{spanning tree} of a connected graph is a subset of $n - 1$ edges that connects all vertices, and a \defn{minimum
spanning tree} (MST) is one of minimum total weight. Every MST algorithm is justified by two properties. A \emph{cut} is a
split of the vertices into two non-empty sides; an edge \emph{crosses} the cut if its ends are on different sides.

\begin{itemize}
  \item \textbf{Cut property.} For any cut, a lightest edge crossing it belongs to some MST.
  \item \textbf{Cycle property.} For any cycle, an edge that is strictly heavier than all the others on the cycle belongs to no
        MST.
\end{itemize}

\emph{Proof of the cut property (an exchange argument).} Let $e$ be a lightest edge crossing a cut, and $T$ an MST without $e$.
Adding $e$ to $T$ creates a cycle, which crosses the cut at $e$ and therefore at least once more, at some edge $f$ of $T$
(Figure~\ref{fig:cutprop}). Replacing $f$ by $e$ gives a spanning tree of weight $w(T) - w(f) + w(e) \le w(T)$, because $e$ is a
lightest crossing edge: it is also an MST, and it contains $e$. The cycle property follows the same way: removing the heaviest
edge of a cycle from any spanning tree and reconnecting with another edge of the cycle never increases the weight.

\begin{figure}[htbp]
\centering
\begin{tikzpicture}
  \fill[fgsky!18, rounded corners=10pt] (-0.6, -1.6) rectangle (2.2, 1.6);
  \fill[fgorange!15, rounded corners=10pt] (3.4, -1.6) rectangle (6.2, 1.6);
  \node[font=\sffamily\scriptsize, text=fgblue] at (0.8, 1.35) {side $S$};
  \node[font=\sffamily\scriptsize, text=fgorange!80!black] at (4.8, 1.35) {side $V \setminus S$};
  \node[vtxs] (a) at (0, 0.6) {}; \node[vtxs] (b) at (1.5, 0.8) {}; \node[vtxs] (c) at (0.6, -0.9) {};
  \node[vtxs] (d) at (4.2, 0.8) {}; \node[vtxs] (e) at (5.5, 0.3) {}; \node[vtxs] (f) at (4.5, -0.9) {};
  \draw[ed] (a) -- (b) (a) -- (c) (d) -- (e) (e) -- (f);
  \draw[ged] (b) -- node[wl, pos=0.25] {$e$: 2} (d);
  \draw[hed] (c) -- node[wl, pos=0.22] {$f$: 5} (f);
  \draw[ed] (b) -- node[wl, pos=0.72] {7} (f);
  \draw[line width=1pt, draw=fgink!60, dashed] (2.8, -1.8) -- (2.8, 1.8);
  \node[notel, anchor=west, text width=5.6cm] at (6.6, 0) {Three edges cross the cut (dashed line). If a spanning tree uses
    $f$ but not the lightest crossing edge $e$, swapping $f$ for $e$ keeps it spanning (the cycle through $e$ crosses the cut
    twice) and does not make it heavier.};
\end{tikzpicture}
\caption{The cut property: a lightest edge crossing any cut can always be part of a minimum spanning tree.}
\label{fig:cutprop}
\end{figure}

\subsection{Kruskal's algorithm}

\defn{Kruskal's algorithm} sorts the edges by weight and adds each edge whose endpoints are not yet connected, which a DSU checks
in $\Oh{\alpha(n)}$ (Figure~\ref{fig:kruskal}). Each accepted edge is a lightest edge crossing the cut between its endpoint's
component and the rest, so it is safe by the cut property; each rejected edge is the heaviest edge of the cycle it would close.
The running time is dominated by sorting: $\Oh{m \log m}$.

\begin{figure}[htbp]
\centering
\newcommand{\krpanel}[3]{%
\begin{tikzpicture}[x=0.95cm, y=0.95cm]
  \coordinate (A) at (0, 1); \coordinate (B) at (1.6, 2); \coordinate (C) at (3.2, 1);
  \coordinate (D) at (0.8, -0.3); \coordinate (E) at (2.4, -0.3);
  \foreach \ea/\eb in {A/B, B/C, A/D, B/D, B/E, C/E, D/E} \draw[ed] (\ea) -- (\eb);
  \foreach \ea/\eb in {#1} \draw[ged] (\ea) -- (\eb);
  \foreach \ea/\eb in {#2} \draw[line width=1.8pt, draw=fgorange] (\ea) -- (\eb);
  \foreach \ea/\eb in {#3} \draw[line width=1.6pt, draw=fgverm, dashed] (\ea) -- (\eb);
  \foreach \ea/\eb/\w in {A/B/2, B/C/6, A/D/3, B/D/5, B/E/4, C/E/1, D/E/7}
    \node[wl, font=\sffamily\tiny] at ($(\ea)!0.5!(\eb)$) {\w};
  \foreach \vn in {A, B, C, D, E} \node[vtxs] at (\vn) {\vn};
\end{tikzpicture}}
\begin{filmstrip}[3]
  \begin{stepbox}{1. C--E (1): accept}
    \krpanel{A/A}{C/E}{A/A}
  \end{stepbox}
  \begin{stepbox}{2. A--B (2): accept}
    \krpanel{C/E}{A/B}{A/A}
  \end{stepbox}
  \begin{stepbox}{3. A--D (3): accept}
    \krpanel{C/E, A/B}{A/D}{A/A}
  \end{stepbox}
  \begin{stepbox}{4. B--E (4): accept}
    \krpanel{C/E, A/B, A/D}{B/E}{A/A}
  \end{stepbox}
  \begin{stepbox}{5. B--D (5): same set, reject}
    \krpanel{C/E, A/B, A/D, B/E}{A/A}{B/D}
  \end{stepbox}
  \begin{stepbox}{6. B--C (6), D--E (7): reject}
    \krpanel{C/E, A/B, A/D, B/E}{A/A}{B/C, D/E}
  \end{stepbox}
\end{filmstrip}
\caption{Kruskal's algorithm. Edges are taken in order of weight; an edge is accepted (orange, then green) when it joins two
different components and rejected (red, dashed) when its endpoints are already connected. The MST has weight
$1 + 2 + 3 + 4 = 10$.}
\label{fig:kruskal}
\end{figure}

\begin{problem}{Road Reparation (CSES)}
Repairing each road has a cost. Find the minimum total cost of repairing roads so that all cities are connected, or report that
it is impossible.
\end{problem}

\begin{code}
void solve() {
    int n, m;
    cin >> n >> m;
    vector<array<ll, 3>> e(m);  // (cost, a, b): sorting orders by cost
    for (auto& [c, a, b] : e) cin >> a >> b >> c;
    sort(all(e));
    DSU d(n + 1);
    ll total = 0;
    int used = 0;
    for (auto [c, a, b] : e)
        if (d.unite(a, b)) total += c, used++;  // skip edges inside a component
    if (used == n - 1) print(total);
    else print("IMPOSSIBLE");
}
\end{code}

\usetemplate{Kruskal MST}{Ctrl+Alt+T, K}{cpkruskal}

\subsection{Prim's algorithm}

\defn{Prim's algorithm} grows a single tree from any start vertex, always adding the cheapest edge that leaves the tree; each
such edge is a lightest edge crossing the cut between the tree and the rest. With a priority queue it runs in $\Oh{m \log n}$,
essentially Dijkstra with the key $w(u, v)$ instead of $d(u) + w(u, v)$. Its real strength is the dense case: when every pair of
$n$ points is an edge ($m \approx n^2/2$), the array version below runs in $\Oh{n^2}$ without ever building the edge list, while
Kruskal would have to sort $n^2/2$ edges.

\begin{problem}{Connecting points}
Connect $n \le 2000$ points in the plane by a network of segments of minimum total cost, where a segment costs the squared
distance between its endpoints.
\end{problem}

\begin{code}
void solve() {
    int n;
    cin >> n;
    vector<pll> p(n);
    for (auto& [x, y] : p) cin >> x >> y;
    auto cost = [&](int i, int j) {  // squared Euclidean distance
        ll dx = p[i].first - p[j].first, dy = p[i].second - p[j].second;
        return dx * dx + dy * dy;
    };
    // O(n^2) Prim: best[v] = cheapest edge from the tree to v
    vector<ll> best(n, LLONG_MAX);
    vector<bool> in(n, false);
    best[0] = 0;
    ll total = 0;
    rep(it, n) {
        int u = -1;  // the cheapest vertex outside the tree
        rep(v, n) if (!in[v] && (u == -1 || best[v] < best[u])) u = v;
        in[u] = true;
        total += best[u];
        rep(v, n) if (!in[v]) chmin(best[v], cost(u, v));
    }
    print(total);
}
\end{code}

\usetemplate{Prim MST (dense graphs)}{Ctrl+Alt+T, Shift+5}{cpprim}

\subsection{Thresholds and the Kruskal reconstruction tree}

The MST also solves \emph{bottleneck} questions. For every pair of vertices, the path between them in an MST minimises the
largest edge weight among all paths (a \emph{minimax} path). So ``the smallest range $D$ such that every cow can talk to every
other, relaying messages'' is the largest edge of the MST (USACO Moocast), and ``can $u$ reach $v$ using only edges of weight at
most $x$?'' depends only on the MST.

For many such queries, the offline recipe is: sort the edges by weight and the queries by threshold, and add edges with a DSU
while sweeping the threshold upwards. At each query, the DSU contains exactly the edges of weight at most the threshold.

\begin{problem}{Reachability under a weight limit}
For each of $q$ queries $(v, x)$, count the vertices reachable from $v$ using only edges of weight at most $x$.
\end{problem}

\begin{code}
void solve() {
    int n, m, q;
    cin >> n >> m >> q;
    vector<array<ll, 3>> e(m);  // (weight, a, b)
    for (auto& [w, a, b] : e) cin >> a >> b >> w;
    sort(all(e));
    vector<array<ll, 3>> qs(q);  // (limit, vertex, index)
    rep(i, q) cin >> qs[i][1] >> qs[i][0], qs[i][2] = i;
    sort(all(qs));
    DSU d(n + 1);
    vi ans(q);
    int j = 0;
    for (auto [x, v, i] : qs) {  // queries by increasing limit
        while (j < m && e[j][0] <= x) d.unite(e[j][1], e[j][2]), j++;
        ans[i] = d.size(v);  // everything connected with edges <= x
    }
    for (ll a : ans) print(a);
}
\end{code}

When the queries must be answered \emph{online}, record Kruskal's merges as a tree instead: every accepted edge of weight $w$
becomes a new node, labelled $w$, whose two children are the roots of the components it joins (Figure~\ref{fig:krt}). This
\defn{Kruskal reconstruction tree} has the original vertices as leaves, and labels increase towards the root. The vertices
reachable from $v$ with edges of weight at most $x$ are exactly the leaves below the highest ancestor of $v$ with label at most
$x$, which binary lifting finds in $\Oh{\log n}$ (Chapter~\ref{ch:trees}); counts, sums and ranks over those leaves become
subtree queries.

\begin{figure}[htbp]
\centering
\begin{tikzpicture}
  \begin{scope}
    \node[vtx] (g1) at (0, 1.2) {1}; \node[vtx] (g2) at (1.6, 1.2) {2};
    \node[vtx] (g3) at (1.6, -0.3) {3}; \node[vtx] (g4) at (0, -0.3) {4};
    \draw[ged] (g1) -- node[wl] {3} (g2); \draw[ged] (g2) -- node[wl] {1} (g3); \draw[ged] (g3) -- node[wl] {4} (g4);
    \draw[ed] (g1) -- node[wl] {5} (g4);
    \node[lbl] at (0.8, -1.1) {graph (MST edges green)};
  \end{scope}
  \begin{scope}[xshift=5.2cm, yshift=1.4cm]
    \node[vtxs, rectangle, rounded corners=2pt, cool] (c) at (0, 0) {4};
    \node[vtxs, rectangle, rounded corners=2pt, cool] (b) at (-0.9, -0.9) {3};
    \node[vtxs, rectangle, rounded corners=2pt, cool] (a) at (-0.3, -1.8) {1};
    \node[vtx] (l1) at (-1.6, -1.8) {1}; \node[vtx] (l2) at (-0.8, -2.7) {2}; \node[vtx] (l3) at (0.2, -2.7) {3};
    \node[vtx] (l4) at (0.9, -0.9) {4};
    \draw[ed] (c) -- (b) (c) -- (l4) (b) -- (l1) (b) -- (a) (a) -- (l2) (a) -- (l3);
    \node[lbl] at (-0.3, -3.5) {Kruskal reconstruction tree};
  \end{scope}
  \node[notel, anchor=west, text width=5.0cm] at (7.4, 0.4) {Merges in order: 2--3 (1), 1--2 (3), 3--4 (4); the edge 1--4 (5)
    is skipped. From leaf 2 with limit $x = 3$, the highest ancestor with label $\le 3$ is the node 3, whose leaves $\{1, 2,
    3\}$ are exactly what 2 reaches.};
\end{tikzpicture}
\caption{A graph and its Kruskal reconstruction tree. Internal nodes (blue squares) are the accepted edges, labelled with their
weights; labels grow towards the root.}
\label{fig:krt}
\end{figure}

\usetemplate{Kruskal reconstruction tree}{Ctrl+Alt+Y, K}{cpkrt}

\begin{trick}[More MST facts worth knowing]
\begin{itemize}
  \item \textbf{Bor\r{u}vka's algorithm}: in each round, every component adds its cheapest outgoing edge; the number of
        components at least halves, so $\Oh{\log n}$ rounds suffice. It shines when the graph is implicit (weights like
        $a_i \oplus a_j$) and ``the cheapest edge leaving a component'' can be found with a data structure such as a binary
        trie (Section~\ref{sec:xor}).
  \item \textbf{Uniqueness}: if all weights are distinct, the MST is unique.
  \item \textbf{Second-best MST}: add one non-tree edge $(u, v)$ and remove the heaviest edge on the tree path between $u$ and
        $v$; the best such swap gives the second-best spanning tree (maximum on paths by binary lifting).
  \item \textbf{Manhattan MST}: for points with the $L_1$ distance, only $\Oh{n}$ candidate edges (the nearest point in each of
        eight octants) are needed, which brings the MST of $10^5$ points within reach.
\end{itemize}
\end{trick}

% ======================================================================================================================
\section[Prüfer codes]{Prüfer codes and counting trees \hfill \ioi}\label{sec:prufer}

Labelled trees on the vertices $1, \ldots, n$ can be encoded as sequences, which turns questions about trees into questions
about sequences. The \defn{Prüfer code} of a tree is built by repeatedly removing the \emph{smallest} leaf and writing down its
neighbour, until two vertices remain (Figure~\ref{fig:prufer}). The result is a sequence of length $n - 2$ over
$\{1, \ldots, n\}$, and every such sequence arises from exactly one tree. Decoding reverses the process: each vertex $v$ that
appears $k$ times in the sequence has degree $k + 1$, so the leaves at every step are known; at step $i$, the smallest current
leaf is attached to the $i$-th entry, which loses one from its remaining degree.

\begin{figure}[htbp]
\centering
\begin{tikzpicture}
  \node[vtx] (v3) at (0, 0) {3}; \node[vtx] (v4) at (1.8, 0) {4};
  \node[vtx, cool] (v1) at (-1.3, 0.9) {1}; \node[vtx, cool] (v2) at (-1.3, -0.9) {2};
  \node[vtx, cool] (v5) at (3.1, 0.9) {5}; \node[vtx] (v6) at (3.1, -0.9) {6};
  \draw[ed] (v1) -- (v3) (v2) -- (v3) (v3) -- (v4) (v4) -- (v5) (v4) -- (v6);
  \foreach \vv/\kk/\ang in {v1/1/135, v2/2/225, v3/3/90, v5/4/45}
    \node[circle, fill=fgverm, text=white, font=\sffamily\tiny\bfseries, inner sep=1pt] at ($(\vv)+(\ang:0.5)$) {\kk};
  \node[notel, anchor=west, text width=6.8cm] at (4.2, 0) {Remove the smallest leaf, write its neighbour:\\
    \textcircled{\scriptsize 1} leaf 1, write 3 \quad \textcircled{\scriptsize 2} leaf 2, write 3\\
    \textcircled{\scriptsize 3} leaf 3 (a leaf now), write 4 \quad \textcircled{\scriptsize 4} leaf 5, write 4\\
    Two vertices (4, 6) remain: the Prüfer code is $3, 3, 4, 4$. Vertices 3 and 4 have degree 3 and appear twice;
    leaves never appear.};
\end{tikzpicture}
\caption{Computing a Prüfer code. The red numbers give the order in which leaves are removed.}
\label{fig:prufer}
\end{figure}

\begin{problem}{Tree from a Prüfer code}
Reconstruct the tree with a given Prüfer code.
\solution Keep the current leaves in an ordered set (a vertex is a leaf when its remaining degree is 1). For each entry $v$ of
the code, the smallest leaf is attached to $v$; remove it, and decrease the remaining degree of $v$, which becomes a leaf when
it reaches 1. The two vertices left at the end form the last edge. With a set this is $\Oh{n \log n}$; the template does it in
$\Oh{n}$ with a moving pointer.
\end{problem}

\begin{code}
void solve() {
    int n;
    cin >> n;
    vi seq(n - 2);
    cin >> seq;
    vi deg(n + 1, 1);  // a vertex appears (degree - 1) times
    for (int v : seq) deg[v]++;
    set<int> leaves;
    rep(v, 1, n + 1) if (deg[v] == 1) leaves.insert(v);
    for (int v : seq) {  // the smallest leaf was removed, attached to v
        int leaf = *leaves.begin();
        leaves.erase(leaves.begin());
        print(leaf, v);
        if (--deg[v] == 1) leaves.insert(v);
    }
    print(*leaves.begin(), *next(leaves.begin()));  // the last edge
}
\end{code}

\usetemplate{Prufer code (tree <-> sequence)}{Ctrl+Alt+U, U}{cpprufer}

The bijection gives two famous counts at once. There are $n^{n-2}$ labelled trees on $n$ vertices (\emph{Cayley's formula}),
because there are $n^{n-2}$ sequences. And since vertex $i$ appears exactly $d_i - 1$ times, the number of labelled trees in which
vertex $i$ has degree $d_i$ is the number of arrangements of that multiset,
\[
  \binom{n - 2}{d_1 - 1, \; d_2 - 1, \; \ldots, \; d_n - 1} = \frac{(n - 2)!}{\prod_i (d_i - 1)!} ,
\]
provided all $d_i \ge 1$ and $\sum_i (d_i - 1) = n - 2$ (otherwise there are none).

\begin{code}
void solve() {
    int n;
    cin >> n;
    vi d(n);
    cin >> d;
    if (n == 1) {  // a single vertex: the only tree has degree 0
        print(d[0] == 0 ? 1 : 0);
        return;
    }
    // Prufer: vertex i appears d_i - 1 times in a sequence of length n - 2
    ll total = 0;
    for (ll x : d) {
        if (x < 1) {
            print(0);
            return;
        }
        total += x - 1;
    }
    if (total != n - 2) {
        print(0);
        return;
    }
    vi fact(n, 1);
    rep(i, 1, n) fact[i] = fact[i - 1] * i % MOD;
    auto power = [](ll b, ll e) {
        ll r = 1;
        for (b %= MOD; e; e >>= 1, b = b * b % MOD)
            if (e & 1) r = r * b % MOD;
        return r;
    };
    ll ans = fact[n - 2];  // (n - 2)! / prod (d_i - 1)!
    for (ll x : d) ans = ans * power(fact[x - 1], MOD - 2) % MOD;
    print(ans);
}
\end{code}

\begin{trick}[Counting spanning structures]
The generalisation of Cayley's formula is the number of spanning forests of the complete graph in which given vertices lie in
different trees: $k\, n^{n - k - 1}$ forests of $n$ labelled vertices with $k$ trees, where $k$ specified vertices are the
roots. For arbitrary graphs, the matrix-tree theorem counts spanning trees as a determinant (Chapter~\ref{ch:combinatorics}).
\end{trick}

% ======================================================================================================================
\section[Topological order]{Topological order \hfill \gold}\label{sec:topo}

A \defn{topological order} of a directed graph lists its vertices so that every edge points forward, from an earlier vertex to a
later one (Figure~\ref{fig:toposort}). It exists exactly when the graph has no directed cycle: a cycle can never point forward all
the way round, and conversely every DAG has a vertex with no incoming edges (follow edges backwards from any vertex; without
cycles the walk must stop), which can go first.

\begin{figure}[htbp]
\centering
\begin{tikzpicture}
  \node[vtx] (a) at (0, 1) {0}; \node[vtx] (b) at (0, -0.4) {1}; \node[vtx] (c) at (1.6, 0.3) {2};
  \node[vtx] (d) at (3.2, 1) {3}; \node[vtx] (e) at (3.2, -0.4) {4}; \node[vtx] (f) at (4.8, 0.3) {5};
  \draw[ded] (a) -- (c); \draw[ded] (b) -- (c); \draw[ded] (c) -- (d); \draw[ded] (c) -- (e); \draw[ded] (d) -- (f);
  \draw[ded] (e) -- (f); \draw[ded] (a) -- (d);
  \begin{scope}[xshift=6.6cm, yshift=0.3cm]
    \foreach \vv [count=\ii] in {0, 1, 2, 3, 4, 5} \node[vtxs, cool] (t\ii) at (\ii*1.1, 0) {\vv};
    \foreach \ia/\ib in {1/3, 2/3, 3/4, 3/5, 4/6, 5/6} \draw[ded] (t\ia) to[bend left=38] (t\ib);
    \draw[ded] (t1) to[bend right=38] (t4);
    \node[lbl] at (3.85, -1.05) {every edge points to the right};
  \end{scope}
\end{tikzpicture}
\caption{A DAG and one of its topological orders. Kahn's algorithm produces it by repeatedly removing a vertex with no incoming
edges: first 0 and 1, then 2, then 3 and 4, then 5.}
\label{fig:toposort}
\end{figure}

\defn{Kahn's algorithm} does exactly that: keep the in-degrees and a queue of vertices whose in-degree is 0; pop a vertex, append it
to the order, and decrease the in-degrees of its successors, pushing those that reach 0. If fewer than $n$ vertices come out, the
remaining ones all have an incoming edge from another remaining vertex, so they contain a cycle. The alternative is DFS: a vertex
finishes only after everything reachable from it, so the vertices in \emph{decreasing} finishing time form a topological order.

\begin{problem}{Course Schedule (CSES)}
Find an order in which to take $n$ courses respecting $m$ prerequisites ``course $a$ before course $b$'', or report that none
exists.
\end{problem}

\begin{code}
void solve() {
    int n, m;
    cin >> n >> m;
    vvi adj(n + 1);
    vi indeg(n + 1, 0);
    rep(i, m) {
        int a, b;
        cin >> a >> b;
        adj[a].push_back(b);
        indeg[b]++;
    }
    queue<int> q;  // courses whose prerequisites are all taken
    rep(v, 1, n + 1) if (!indeg[v]) q.push(v);
    vi order;
    while (!q.empty()) {
        int u = q.front();
        q.pop();
        order.push_back(u);
        for (int v : adj[u])
            if (--indeg[v] == 0) q.push(v);
    }
    if (sz(order) < n) print("IMPOSSIBLE");  // the rest lie on cycles
    else print(order);
}
\end{code}

\usetemplate{Topological sort (Kahn)}{Ctrl+Alt+T, O}{cptopo}

\begin{recognize}
Dependencies (``task $a$ before task $b$''), prerequisites, ``is there an order consistent with these comparisons?'', build
systems, and every DP over a DAG. The lexicographically smallest topological order uses a min-heap instead of the queue; the
order that places a particular vertex as early as possible is found by running Kahn on the \emph{reversed} graph and taking the
largest available vertex each time.
\end{recognize}

\subsection{Dynamic programming along a topological order}

In topological order, every vertex is processed after all its predecessors, so any quantity defined through incoming edges is
computed in one pass: longest and shortest paths (Section~\ref{sec:dagpaths}), the number of paths, reachability with extra
conditions (Figure~\ref{fig:gameroutes}).

\begin{problem}{Game Routes (CSES)}
Count the routes from level 1 to level $n$ in a game whose teleporters form a DAG, modulo $10^9 + 7$.
\solution The number of routes to $v$ is the sum of the numbers of routes to its predecessors: $\mathit{ways}(v) =
\sum_{u \to v} \mathit{ways}(u)$, with $\mathit{ways}(1) = 1$, evaluated in topological order.
\end{problem}

\begin{figure}[htbp]
\centering
\begin{tikzpicture}[x=1.6cm]
  \foreach \vn/\xx/\yy/\val in {1/0/0/1, 2/1/0.8/1, 3/1/-0.8/2, 4/2/0.8/3, 5/2/-0.8/2, 6/3/0/5} {
    \node[vtx] (g\vn) at (\xx, \yy) {\vn};
    \node[font=\sffamily\scriptsize\bfseries, text=fgblue, above=1pt] at (g\vn.north) {\val};
  }
  \draw[ded] (g1) -- (g2); \draw[ded] (g1) -- (g3); \draw[ded] (g2) -- (g3); \draw[ded] (g2) -- (g4);
  \draw[ded] (g3) -- (g4); \draw[ded] (g3) -- (g5); \draw[ded] (g4) -- (g6); \draw[ded] (g5) -- (g6);
  \node[notel, anchor=west, text width=5.4cm] at (3.5, 0) {Blue: the number of routes from 1. Each value is the sum over the
    incoming edges: $\mathit{ways}(4) = \mathit{ways}(2) + \mathit{ways}(3) = 1 + 2$, and $\mathit{ways}(6) = 3 + 2 = 5$.};
\end{tikzpicture}
\caption{Counting paths in a DAG by summing over incoming edges in topological order.}
\label{fig:gameroutes}
\end{figure}

\begin{code}
void solve() {
    int n, m;
    cin >> n >> m;
    vvi adj(n + 1);
    vi indeg(n + 1, 0);
    rep(i, m) {
        int a, b;
        cin >> a >> b;
        adj[a].push_back(b);
        indeg[b]++;
    }
    vi order, st;  // a topological order by Kahn's algorithm
    rep(v, 1, n + 1) if (!indeg[v]) st.push_back(v);
    while (!st.empty()) {
        int u = st.back();
        st.pop_back();
        order.push_back(u);
        for (int v : adj[u])
            if (--indeg[v] == 0) st.push_back(v);
    }
    vi ways(n + 1, 0);  // ways[v]: routes from 1 to v, modulo MOD
    ways[1] = 1;
    for (int u : order)
        for (int v : adj[u]) ways[v] = (ways[v] + ways[u]) % MOD;
    print(ways[n]);
}
\end{code}

\begin{trick}[Cycles? Condense first]
If the graph has cycles, contract every strongly connected component to a single vertex (Chapter~\ref{ch:connectivity}). The
result is a DAG, and many questions about walks that may revisit vertices (the most coins collectable, the largest set of
mutually reachable rooms) become DP on that DAG (CSES Coin Collector).
\end{trick}

\begin{summary}[Chapter summary]
\begin{itemize}
  \item A DSU with union by size and path compression costs $\Oh{\alpha(n)}$ per operation; offsets along the tree maintain
        differences, parities and other group-valued relations.
  \item The cut property (a lightest crossing edge is safe) justifies Kruskal ($\Oh{m \log m}$, with a DSU) and Prim
        ($\Oh{n^2}$ on dense graphs, $\Oh{m \log n}$ with a heap).
  \item MST paths are minimax paths; threshold queries are answered offline by sorting, or online with a Kruskal reconstruction
        tree.
  \item Prüfer codes are a bijection between labelled trees and sequences of length $n - 2$: $n^{n-2}$ trees, and multinomial
        counts for given degrees.
  \item A graph has a topological order if and only if it is acyclic; Kahn's algorithm finds one or detects the cycle, and DP
        along it counts paths and optimises over them.
\end{itemize}
\end{summary}

\begin{practice}
\ppgroup{Union--find and spanning trees}
\begin{pplist}
  \item Road Reparation
  \item Road Construction
  \item Network Breakdown
  \item New Roads Queries
\end{pplist}
\ppgroup{Topological order and DAGs}
\begin{pplist}
  \item Course Schedule
  \item Longest Flight Route
  \item Game Routes
  \item Course Schedule II
  \item Coin Collector
\end{pplist}
\end{practice}
